% Generated draft: requires mathematical and source review.
\documentclass[11pt]{article}
\usepackage[margin=25mm]{geometry}
\usepackage{fontspec}
\setmainfont{Latin Modern Roman}
\usepackage{amsmath,amssymb,mathtools,unicode-math}
\setmathfont{Latin Modern Math}
\usepackage{enumitem,needspace,xurl,hyperref}
\setlength{\parindent}{0pt}
\setlength{\parskip}{5pt}
\setlength{\emergencystretch}{3em}
\newcommand{\N}{\mathbb{N}}
\newcommand{\Z}{\mathbb{Z}}
\newcommand{\Q}{\mathbb{Q}}
\newcommand{\R}{\mathbb{R}}
\newcommand{\C}{\mathbb{C}}
\newcommand{\F}{\mathbb{F}}
\newcommand{\A}{\mathbb{A}}
\newcommand{\PP}{\mathbb{P}}
\newcommand{\E}{\mathbb{E}}
\newcommand{\eps}{\varepsilon}
\newcommand{\dd}{\,\mathrm{d}}
\begin{document}
\section*{UnsolvedMath ID: 1101201}
\textbf{Source problem code:} catalogue entry
\subsection{Definitions and mathematical statement}
Let $S$ be a compact non-orientable hyperbolic surface. A measured geodesic lamination is a closed union of pairwise disjoint complete geodesics equipped with locally finite transverse measures invariant under sliding a transversal along leaves. A leaf is two-sided when its normal line bundle in $S$ is trivial; for a closed leaf this means that it has an annular, rather than a M\"obius-band, neighborhood. Let $\mathcal{ML}^{+}(S)$ denote the measured laminations all of whose leaves are two-sided. Give measured laminations their usual topology of convergence of transverse measures. A simplicial measured lamination is a finite weighted union
\[
 \lambda_m=\sum_{j=1}^{r_m}a_{m,j}\gamma_{m,j},\qquad a_{m,j}>0,
\]
where the $\gamma_{m,j}$ are pairwise disjoint simple closed geodesics; the empty union represents the zero lamination. It has two-sided leaves when every $\gamma_{m,j}$ is two-sided.

\textbf{Question attributed to Levitt.} For every such $S$ and every $\lambda\in\mathcal{ML}^{+}(S)$, is there a sequence of simplicial measured laminations $\lambda_m\in\mathcal{ML}^{+}(S)$ converging to $\lambda$? Equivalently, is the subset of two-sided weighted multicurves dense in $\mathcal{ML}^{+}(S)$, for the topology inherited from measured laminations?

The record marks the question solved. Its research summary claims the stronger approximation by single weighted two-sided simple closed geodesics, namely $a_m\gamma_m\to\lambda$, for compact non-orientable hyperbolic surfaces of \emph{non-exceptional type}. It attributes this to Erlandsson--Gendulphe--Pasquinelli--Souto, Theorem 1.2 together with Lemma 2.3 (GAFA, 2023), and reports a partial version for, for example, orientable ergodic laminations by Khan (2021/2023). The meaning and extent of ``non-exceptional type'' and the boundary convention are absent from the supplied record; consequently the claimed result is not silently extended to every surface in the literal question.
\subsection{Short English statement}
Levitt asks whether every measured lamination consisting of two-sided leaves on a compact non-orientable hyperbolic surface is a limit of finite weighted unions of two-sided simple closed geodesics. The supplied summary reports even approximation by single weighted curves, but only for an unspecified class of non-exceptional surfaces.
\subsection{Sources}
\begin{itemize}
\item Local UnsolvedMath registry, permanent ID 1101201, source code catalogue entry, ``Questions in Geometric Group Theory --- Q 12.1''; attribution: Levitt. Supplied status: solved; research classification: SOLVED-IN-LITERATURE; updated 2026-08-06.
\item Mladen Bestvina, \emph{Questions in Geometric Group Theory}, source list dated 2004, Question 12.1, PDF page 19, \url{https://www.math.utah.edu/~bestvina/eprints/questions-updated.pdf}. The record reports access on 2026-07-29 and manual transcription. Its background says that the rendered author PDF presents an open item and current status was not independently verified; status and rights require review, with the PDF described as publicly author-hosted.
\item Supplied research summary: Erlandsson--Gendulphe--Pasquinelli--Souto, GAFA (2023), Theorem 1.2 and Lemma 2.3, reportedly citing Bestvina's question; and Khan (2021/2023), reportedly giving partial progress. No article titles, complete bibliographic entries, or URLs for these papers are supplied.
\end{itemize}
\end{document}
